%% introduction_ko.tex
\chapter{서론}
\label{chap:introduction_ko}

\section{연구 배경 및 필요성}
\label{sec:intro_motivation_ko}

현대 기업은 API 기술 참조 문서, 인사 정책 매뉴얼, 시스템 아키텍처 설계서, 회의록, 보안 컴플라이언스 가이드, 출장 비용 정책 등 매우 다양한 형태의 지식 문서를 대규모로 생성하고 관리한다. 이러한 이종 문서 코퍼스의 규모가 지속적으로 증가함에 따라, 필요한 정보를 신속하고 정확하게 검색하는 문제는 조직 운영의 핵심 과제로 자리매김하였다. 수백 개의 정책 문서를 검색하는 직원, 시스템 설계 아카이브를 조회하는 엔지니어, 컴플라이언스 검사를 수행하는 자동화 에이전트 등 다양한 사용자와 시스템이 문서 검색에 의존하고 있다.

검색 증강 생성(Retrieval-Augmented Generation, RAG)은 지식 집약적 자연어 처리 과제를 위한 선도적 패러다임으로 부상하였다 \citep{lewis2020rag}. RAG 시스템은 코퍼스에서 후보 문서나 구절을 검색한 후, 검색된 구절에 조건화하여 생성 언어 모델의 출력을 생성한다. 따라서 검색 단계는 시스템 전체의 핵심적인 선행 요소이다. 관련 문서가 검색되지 않으면 어떤 수준의 생성 품질로도 이를 보완할 수 없다. 검색 품질은 \emph{커버리지}(관련 문서를 모두 찾는가?)와 \emph{순위 품질}(가장 관련성 높은 문서를 상위에 배치하는가?)이라는 두 가지 차원에서 평가된다.

기업 코퍼스는 표준적인 검색 방법에 독특한 도전을 제기한다. 웹 검색이나 개방형 질의응답과 달리, 기업 질의는 단순한 단일 홉(single-hop) 사실 조회부터 여러 문서에 걸친 복잡한 다중 홉(multi-hop) 추론, 시간적 순서 질의, 정책 예외 조항 추출, 다속성 집계 질의까지 매우 다양한 추론 유형을 포괄한다. 각 추론 유형은 서로 다른 검색 신호로부터 이익을 얻는다. 시간적 질의는 밀집 임베딩이 포착하지 못하는 정확한 날짜 토큰 매칭이 필요하고, 예외 조항 질의는 부정 및 조건 언어의 의미적 식별이 요구되며, 다중 홉 질의는 광범위한 문서 커버리지가 효과적이다. 따라서 어떠한 단일 검색 방법도 이 모든 기업 질의 유형에 최적화될 수 없다.

\emph{적응형 검색 계획(Adaptive Retrieval Planning)}은 모든 질의에 단일 고정 전략을 적용하는 대신, 입력 질의의 특성에 가장 적합한 검색 전략을 동적으로 선택함으로써 이 문제를 해결한다. 검색 계획기는 질의의 추론 유형, 의미적 내용, 개체명, 난이도 요인 등을 분석하여 가장 적합한 검색 백엔드로 질의를 라우팅한다. 이 라우팅 결정이 각 질의 유형에 대한 실증적 최적 전략과 일치한다면, 계획기는 최소한의 추가 계산 비용으로 검색 품질을 향상시킬 수 있다.

본 논문은 기업 지식 관리 맥락에서 적응형 검색 계획을 연구한다. 핵심 질문은 다음과 같다: 검색 계획기 — 경량 결정론적 규칙 시스템이든 대규모 언어 모델이든 — 가 최선의 고정 전략 검색 기준선을 능가할 수 있는가, 그리고 그 비용은 어느 정도인가?

\section{연구 문제}
\label{sec:intro_problem_ko}

본 논문은 다음과 같은 핵심 연구 문제를 다룬다:

\begin{quote}
\emph{다양한 유형의 문서와 다양한 추론 유형에 걸친 질의가 존재하는 이종 기업 지식 코퍼스에서, 가장 효과적이고 비용 효율적인 검색 전략 선택 메커니즘은 무엇인가? 그리고 적응형 검색 계획은 고정 전략 검색 기준선에 비해 어느 정도 성능을 개선하는가?}
\end{quote}

이 연구 문제는 세 가지 하위 문제로 분해된다:

\begin{enumerate}
    \item \textbf{검색 기준선 선택}: BM25(희소), 밀집 임베딩 검색, Hybrid RRF 중 전체 성능, 문서 범주별, 질의 추론 유형별로 어떤 방법이 가장 우수한가?
    \item \textbf{규칙 기반 적응형 계획}: 질의 메타데이터를 검색 전략으로 매핑하는 소규모 결정론적 규칙 집합이 최선의 고정 전략 기준선을 능가할 수 있는가?
    \item \textbf{LLM 기반 적응형 계획}: GPT-5와 같은 대규모 언어 모델이 고정 전략 기준선과 규칙 기반 계획기를 모두 능가하는 제로샷 검색 전략 계획을 수행할 수 있으며, 그 품질 향상은 비용 대비 정당화되는가?
\end{enumerate}

\section{연구 질문}
\label{sec:intro_rqs_ko}

본 연구는 7개의 연구 질문(RQ1--RQ7)을 중심으로 탐구한다:

\begin{description}
    \item[RQ1] 하이브리드 검색(BM25 + Dense)은 기업 지식 질의에서 단일 방법 기준선을 능가하는가?
    \item[RQ2] 다양한 기업 문서 범주에서 어떤 검색 방법이 가장 우수한가?
    \item[RQ3] 질의 추론 유형에 따라 검색 성능은 어떻게 달라지는가?
    \item[RQ4] 규칙 기반 적응형 계획기가 고정 전략 검색을 능가할 수 있는가?
    \item[RQ5] LLM 기반 검색 계획에 최적인 프롬프트 전략은 무엇인가?
    \item[RQ6] GPT-5는 전략 선택 정확도(SSA)에서 규칙 기반 계획기를 능가하는가?
    \item[RQ7] LLM 기반 검색 계획의 주요 실패 유형은 무엇인가?
\end{description}

RQ1--RQ3은 검색 기준선의 특성을 확립하고, RQ4는 결정론적 적응형 계획을 평가하며, RQ5--RQ7은 LLM 기반 계획을 평가한다. 이 7개의 연구 질문은 기업 지식 관리에서 검색 전략 선택의 전체 설계 공간을 포괄한다.

\subsection*{확장: V2 에이전틱 RAG (RQ8--RQ13)}
\label{sec:intro_rqs_v2_ko}

V1 실험은 모든 정적 계획 접근법이 공유하는 구조적 한계를 드러낸다: 검색 전략 선택이 질의당 한 번 이루어지며, 초기 전략이 실패할 때 수정될 수 없다. V1에서 일관되게 다루기 어려운 것으로 입증된 두 가지 질의 유형 — 예외 처리(Recall@10 $\leq 0.412$)와 시간적 추론(Recall@10 $\leq 0.200$) — 은 근본적으로 다른 아키텍처를 동기화한다. 제~\ref{chap:v2_methodology_ko}--\ref{chap:v2_results_ko}장은 검증 신호에 조건화된 반복적 재검색을 지원하는 LangGraph 상태 기계로 V1 검색 인프라를 래핑하는 V2 에이전틱 RAG 파이프라인을 소개한다. 이 확장을 위한 추가 연구 질문 6개:

\begin{description}
    \item[RQ8]  V2 에이전틱 파이프라인은 최우수 V1 계획기(규칙 계획기, Recall@10$={0.755}$)보다 전체 검색 품질에서 우수한가?
    \item[RQ9]  반복 재검색 루프는 단일 패스 에이전틱 기준선보다 어떤 기여를 하는가?
    \item[RQ10] 신호 매칭(적응형) 검증 스코어링은 순위 순서 스코어링보다 더 나은 검색 검증을 가능하게 하는가?
    \item[RQ11] Oracle 질의 분석이 휴리스틱 LLM 기반 질의 분석에 비해 V2 성능에 얼마나 기여하는가?
    \item[RQ12] 타겟 재검색 메커니즘이 V1에서 식별된 예외 처리 및 시간적 질의 실패 모드를 구체적으로 해결하는가?
    \item[RQ13] 모든 재검색 루프가 소진된 후 V2 에이전틱 파이프라인의 주요 실패 모드는 무엇인가?
\end{description}

\section{연구 기여}
\label{sec:intro_contributions_ko}

본 논문의 기여는 다음 네 가지 측면에서 이루어진다.

\textbf{방법론적 기여}: 합성 기업 지식 데이터셋(SEKD)을 구축하였다. SEKD는 6개 범주에 걸쳐 120개 문서, 1,206개 섹션 인식 청크, 정답 청크 ID 및 oracle 전략 레이블이 포함된 405개 주석 질의로 구성된 완전한 주석 기업 코퍼스이다.

\textbf{실증적 기여}: BM25, 밀집 검색, Hybrid RRF, 규칙 기반 계획기, LLM 기반 계획기를 동일한 통제된 기업 지식 코퍼스에서 비교하는 최초의 체계적 연구를 수행하였다. 각 아키텍처 계층의 성능 기여를 정량화하고, 각 방법이 우수하거나 실패하는 구체적인 질의 유형을 식별하였다.

\textbf{실용적 기여}: 7개 규칙으로 구성된 결정론적 계획기가 \$0의 운영 비용으로 V1 시스템 중 최고의 Recall@10(0.755)을 달성함을 입증하였다. V2 에이전틱 확장은 모든 V1 접근법에서 다루기 어렵던 예외 처리 및 시간적 질의 유형에 대해 반복적 수정 메커니즘을 제공한다.

\textbf{분석적 기여}: LLM 기반 검색 계획의 주요 실패 유형 — 세밀한 단일 홉 질의에서의 범주 혼동(오류의 55.3\%)과 시간적 질의의 체계적 오분류(SSA 0\%) — 을 규명하고, 모든 평가된 검색 방법에 걸친 예외 처리 및 시간적 질의 실패의 구조적 원인을 분석하였다. V2 파이프라인의 실패 모드 분류법은 이 분석을 에이전틱 설정으로 확장한다.

\textbf{시스템 기여}: V2 에이전틱 RAG 시스템을 5개의 전문화된 에이전트와 구조화된 평가 하네스를 포함하는 완전히 운영 가능한 LangGraph 파이프라인으로 구현하였다. 이 구현은 향후 반복적 기업 검색 연구를 위한 재현 가능한 기반 역할을 한다.

\section{논문 구성}
\label{sec:intro_structure_ko}

본 논문의 나머지 부분은 다음과 같이 구성된다. 제2장은 정보 검색, 밀집 검색, 하이브리드 융합 방법, LLM 증강 시스템에 관한 관련 연구를 검토한다. 제3장은 SEKD 코퍼스 구축, 검색 시스템 설계, 계획기 아키텍처, 평가 프레임워크를 기술한다. 제4장은 V1 시스템의 실험 결과를 제시한다. 제5장은 각 V1 시스템의 성능 원인을 분석하고 기업 RAG 설계에 대한 함의를 도출한다. 제6장은 V1 주요 연구 발견을 요약한다. 제7장은 V2 에이전틱 RAG 파이프라인 아키텍처와 평가 설계를 소개한다. 제8장은 RQ8--RQ13을 다루는 V2 평가 결과를 제시한다.
